\par\noindent\textbf{(1)} 均匀背景 \(\psi_0=Ae^{igA^2t}\) 满足NLS，因为 \(i\partial_t\psi_0=-gA^2\psi_0\)，与非线性项相消。设 \(\psi=e^{igA^2t}[A+u(x,t)+iv(x,t)]\)，其中 \(u,v\) 为实小扰动。
\par\noindent\textbf{(2)} 保留一阶扰动得 \(u_t=-(D/2)v_{xx}\)、\(v_t=(D/2)u_{xx}+2gA^2u\)。令扰动 \(\propto e^{i(Kx-\Omega t)}\)，消去振幅得 \(\Omega^2=(D^2K^4/4)-DgA^2K^2\)。这个色散由两个实变量的线性矩阵行列式得到。
\par\noindent\textbf{(3)} 若 \(Dg>0\)，当 \(0<K^2<4gA^2/D\)（按 \(g/D>0\)）时 \(\Omega^2<0\)，振幅按 \(e^{\Gamma t}\) 增长，\(\Gamma=|K|\sqrt{DgA^2-D^2K^2/4}\)。对 \(K^2\) 求极值得 \(K_*^2=2gA^2/D\)，最大增长率 \(\Gamma_{\max}=|g|A^2\)。
\par\noindent\textbf{(4)} 若 \(Dg<0\)，色散 \(\Omega^2>0\) 对所有非零 \(K\) 成立，平面波线性谱稳定；但这不证明大振幅非线性摄动绝对稳定。有限带宽、耗散和更高阶色散会改变不稳定带宽与阈值。
